\section{Thiết kế thực nghiệm}
\label{sec:experimental-methodology}

\subsection{Câu hỏi nghiên cứu}
\label{sec:rqs}

\begin{itemize}
  \item \textbf{RQ1 (duy trì chất lượng theo thang LLM-as-judge).} Khi model nền
  nhỏ dần trên ladder Qwen 0.8B--9B, AlignmentGraph-ISD có duy trì được chất lượng
  thiết kế đo bằng LLM-as-judge của ISD-Agent-Bench so với sáu benchmark agent
  không: hệ đứng ở đâu tại từng kích thước, và khoảng cách có nới rộng khi model
  nhỏ đi không?
  \item \textbf{RQ2 (duy trì chất lượng theo thang tín hiệu alignment tường
  minh).} Khi được làm tường minh bằng giao thức đo không dùng LLM, các tín hiệu
  alignment nội tại của thiết kế do máy sinh thay đổi thế nào theo kích thước
  model, và khoảng cách so sánh có nới rộng khi model nhỏ đi không?
\end{itemize}
Hai vế của RQ1 --- mức tại từng kích thước và scale interaction --- là đồng hạng
và được kiểm bằng hai phép phân tích riêng (per-size và
difference-in-differences, \Cref{sec:stats}): significance ở một kích thước không
được hiểu nhầm là bằng chứng về quá trình thu nhỏ model. RQ2 chủ ý dừng ở các tín
hiệu đo được; construct validation đòi hỏi bằng chứng chuyên gia ngoài điểm số tự
động.

\subsection{Các điều kiện so sánh}
\label{sec:conditions}

So sánh chính gồm bảy hệ: AlignmentGraph-ISD và sáu benchmark agent (Baseline,
EduPlanner, ADDIE-Agent, RPISD-Agent, Dick-Carey-Agent, ReAct-ISD) ở mọi kích
thước và mọi run.

Cho ablation study, ba arm của hệ đề xuất chạy kèm, mỗi arm tắt đúng một phần cơ
chế trong \Cref{sec:method}: $-$Verifier tắt tầng kiểm tra--repair
(\Cref{sec:qc}) nhưng giữ phân rã và graph-context; $-$GraphContext tắt các lượt
đọc context slice $C(s)$ nhưng giữ nguyên lượt ghi theo contract $W(s)$
(\Cref{sec:alignment-graph}); skeleton tắt cả hai, chỉ giữ specialist
decomposition. Thiết kế này kiểm hai cơ chế xây trên nền decomposition, không
kiểm bản thân decomposition --- nền kiến trúc được kế thừa
(\Cref{sec:decomposition}). Cả ba arm chạy trong cùng các session của ladder, nên
delta arm-trừ-A0 dùng chung judge session và tránh judge drift; A0 trong session
cũng là consistency check với ba run A0 chính.

\subsection{Dữ liệu và ladder model}
\label{sec:dataset}
\label{sec:setup}

Mọi điều kiện chạy trên split \texttt{test\_90} của
ISD-Agent-Bench~\cite{isdagentbench}: 90 scenario được chọn bằng stratified
sampling với seed 42, dùng chung cho mọi hệ. Ladder gồm các checkpoint Qwen
0.8B, 2B, 4B và 9B; nhờ dùng cùng một họ model, chỉ có năng lực thay đổi giữa các
kích thước. Mỗi ô agent--kích thước có ba run độc lập trên toàn bộ 90 scenario.
Scenario và prompt của cả bảy agent đều bằng tiếng Anh theo thiết kế; can thiệp
duy nhất dành riêng cho baseline là dịch prompt scaffolding của ReAct-ISD sang
tiếng Anh mà không thay đổi phương pháp. Manifest ghi checkpoint, prompt, seed,
decoding setting và judge.

\subsection{Deployment envelope và điều kiện serving}
\label{sec:assumptions}

Deployment target là một GPU 12--16\,GB. Với hai byte mỗi parameter, trọng số
bf16 cần xấp xỉ 1.6\,GB cho 0.8B, 4\,GB cho 2B, 8\,GB cho 4B và 18\,GB cho 9B;
model 9B vì vậy không vừa, và 4B là checkpoint lớn nhất trong target. Phần VRAM
còn lại phải chứa KV cache và runtime overhead, nên ngân sách còn khoảng 8k token
mỗi lời gọi. Chuỗi suy luận này kiểm chứng được độc lập với thời điểm thực
nghiệm; 9B là mốc tham chiếu ngoài target, không phải cấu hình triển khai được đề
xuất.

Mọi agent chạy qua vLLM với streaming được bật, thinking bị tắt và generation cap
chung. Chúng tôi giới hạn generation ở 8192 token mỗi lời gọi --- ngân sách mà model 4B để lại trên một GPU 12--16\,GB, là deployment target mà nghiên cứu này tuyên bố. Các pod 24--32\,GB chỉ là tiện nghi thuê máy để chạy ladder song song, không phải envelope được tuyên bố.
Do đó, việc một kiến trúc có hoàn thành trong giới hạn này hay không là deployment
property được đo.

\subsection{Chỉ số đánh giá}
\label{sec:metrics}

Với RQ1, hai judge cố định thuộc các họ tách biệt với agent ---
\texttt{openai/gpt-4o-mini} và \texttt{google/gemini-2.5-flash-lite} --- chấm mọi
artifact; sự tách biệt này giảm nguy cơ
self-preference~\cite{panickssery2024llm}. Điểm scenario là trung bình của hai
judge. Ba tín hiệu judge được báo song song với cách phân tích giống nhau và vị
thế ngang nhau (\Cref{tab:judge-signals}); ADDIE dẫn mạch trình bày vì rubric của
nó khớp construct mà RQ1 hỏi.

\begin{table}[t]
  \caption{Ba tín hiệu judge của RQ1 (thang /100).}
  \label{tab:judge-signals}
  \small
  \begin{tabular}{@{}lp{0.35\columnwidth}p{0.24\columnwidth}@{}}
    \toprule
    Tín hiệu & Construct đo & Vai trò \\
    \midrule
    \texttt{addie\_median} & chất lượng nội dung thiết kế (artifact rubric) &
    dẫn mạch trình bày \\
    \texttt{trajectory\_score} & tuân thủ quy trình tool-use (BFCL) &
    báo song song \\
    \texttt{total\_score} & composite $0.7\,$ADDIE$+0.3\,$Traj của benchmark &
    đối chiếu với benchmark \\
    \bottomrule
  \end{tabular}
\end{table}

Với RQ2, giao thức trong \Cref{sec:protocol} chạy offline trên output đã lưu;
công cụ đo được cấu hình theo preset (\Cref{tab:rq2-arms}): cấu hình chính ghép
encoder Nemotron với checkpoint BERT fine-tuned, và bốn arm sensitivity tạo thành một
star quanh cấu hình chính --- mỗi arm lệch đúng một trục, không phải grid.
Checkpoint BERT luôn trả argmax còn lexicon có thể abstain, nên denominator của
Family B có thể khác nhau giữa hai arm Bloom. Quy tắc đọc panel pre-specified: có
lợi thế tại một kích thước khi ít nhất bốn trong sáu tín hiệu có hướng dương có ý
nghĩa sau hiệu chỉnh Holm trong-tín-hiệu và không tín hiệu nào âm có ý nghĩa; mẫu
ngược lại kết luận bất lợi, còn lại là mixed. Tín hiệu thứ bảy không định hướng
vẫn được báo cáo. Về vận hành: quy trình chấm không dùng LLM và không có thành phần ngẫu nhiên; kết quả tái lập được nhờ pin phiên bản encoder và embedding cache.
Cost telemetry chỉ đọc \texttt{metadata.token\_usage}.

\begin{table}[t]
  \caption{Cấu hình công cụ đo: cấu hình chính và bốn arm sensitivity (star ---
  mỗi arm lệch đúng một trục so với cấu hình chính).}
  \label{tab:rq2-arms}
  \small
  \begin{tabular}{@{}lllp{0.2\columnwidth}@{}}
    \toprule
    Arm & Embedding encoder & Bloom classifier & Trục kiểm tra \\
    \midrule
    Chính & Nemotron-8B~\cite{babakhin2025nemotron} &
    fine-tuned BERT~\cite{li2022bloomclassification} & --- \\
    BGE-M3 & BGE-M3~\cite{chen2024m3embedding} & fine-tuned BERT & lineage encoder \\
    Granite & Granite-311M~\cite{awasthy2026granite} & fine-tuned BERT & lineage encoder \\
    TF--IDF & TF--IDF char $n$-gram & fine-tuned BERT & floor không semantic \\
    Lexicon & Nemotron-8B & verb lexicon~\cite{anderson2001taxonomy} &
    Bloom classifier \\
    \bottomrule
  \end{tabular}
\end{table}

\subsection{Phân tích thống kê}
\label{sec:stats}

Đơn vị suy luận là trung bình của một scenario qua ba run. Với vế mức của RQ1,
paired Wilcoxon signed-rank test so sánh A0 với từng baseline ở mọi kích thước và
trên cả ba tín hiệu judge; RQ2 áp dụng cùng phép kiểm cho từng tín hiệu panel.
Khi paired sample có $n\leq25$, chúng tôi dùng exact permutation null. Chúng tôi
báo paired mean difference, rank-biserial effect và confidence interval bootstrap
95\% từ 10{,}000 lần resample scenario với seed 42.

Hiệu chỉnh Holm bao phủ sáu so sánh proposed-versus-baseline mỗi kích thước,
riêng cho từng tín hiệu. RQ2 còn báo hiệu chỉnh Holm toàn panel trên bảy tín hiệu
và sáu so sánh như multiplicity sensitivity; hiệu chỉnh này không tạo hierarchy
giữa các tín hiệu. Với vế scale của RQ1, gọi $\Delta_s$ là khoảng cách ghép cặp
A0-trừ-baseline tại kích thước $s$. Interaction pre-specified là
$\Delta_{\mathrm{largest}}-\Delta_{\mathrm{smallest}}$, được kiểm tra bằng paired
Wilcoxon cùng bootstrap interval; giá trị âm nghĩa là lợi thế A0 lớn hơn ở đầu
nhỏ. RQ2 dùng cùng interaction, với $\sigma_{\mathrm{meas}}$
dẫn mạch scale và các tín hiệu khác được báo song song.

Với output bị thiếu, layer chính gán \texttt{failure=0}: agent không sinh được
output nghĩa là chất lượng bằng 0, không phải dữ liệu khuyết. Suy luận
complete-case chỉ là sensitivity và không bao giờ là unconditional, vì nó có điều
kiện trên một biến post-treatment: loại bỏ đúng các failure nặng nhất, đặc biệt ở
kích thước nhỏ. RQ1 báo thêm \texttt{failure=min}; RQ2 không dùng layer này vì
tín hiệu đã có floor 0. Output-failure rate được so sánh riêng bằng Fisher exact
test với hiệu chỉnh Holm.
